\subsection{Adjunctions}
\label{subsec:mates-adjunctions}

By \cref{thm:walking-adjunction}, $\func(\Adj,\E)$ is the $(\infty,2)$-category of adjunctions in $\E$, and extending a left adjoint to an adjunction is a property. We first describe $\func(\Adj,\E)$ inside $\func(\Delta^{1},\E)$: this is the mate lemma. We then show that, when $\E$ has lax pullbacks, adjunctions and their morphisms can be described in terms of lax pullbacks, without reference to $\Adj$.

\subsubsection*{The mate lemma}

\begin{lem}[mates]
\label{lem:mates}
Let $\E$ be an $(\infty,2)$-category.
\begin{enumerate}
\item $\ell^{*}\colon \func(\Adj,\E) \to \func(\Delta^{1},\E)$ is a locally full inclusion. It exhibits $\func(\Adj,\E)$ as the locally full sub-$(\infty,2)$-category whose objects are the left adjoints $u\colon x \to y$, and whose $1$-morphisms from $u$ to $u'$ are the commutative squares for which the lax square
\[
\begin{tikzcd}
x \ar[r, "f_{0}"] \ar[d, "u"'] & x' \ar[d, "u'"] \ar[dl, Rightarrow, shorten=3mm, "\varphi"'] \\
y \ar[r, "f_{1}"'] & y'
\end{tikzcd}
\qquad \varphi\colon u'f_{0} \simeq f_{1}u \ ,
\]
is right Beck--Chevalley: $\rmate{\varphi}\colon f_{0}u^{R} \Rightarrow u'^{R}f_{1}$ is invertible.
\item Dually, $\rho^{*}\colon \func(\Adj,\E) \to \func(\Delta^{1},\E)$ is a locally full inclusion. It exhibits $\func(\Adj,\E)$ as the locally full sub-$(\infty,2)$-category whose objects are the right adjoints $u\colon x \to y$, and whose $1$-morphisms from $u$ to $u'$ are the commutative squares for which the lax square
\[
\begin{tikzcd}
x \ar[r, "u"] \ar[d, "f_{0}"'] & y \ar[d, "f_{1}"] \ar[dl, Rightarrow, shorten=3mm, "\psi"'] \\
x' \ar[r, "u'"'] & y'
\end{tikzcd}
\qquad \psi\colon f_{1}u \simeq u'f_{0} \ ,
\]
is left Beck--Chevalley: $\lmate{\psi}\colon u'^{L}f_{1} \Rightarrow f_{0}u^{L}$ is invertible.
\end{enumerate}
\end{lem}

\begin{proof}
We prove (1); (2) is (1) for $\E^{\mathrm{co}}$, since $\Adj^{\mathrm{co}} \simeq \Adj$ exchanges $\ell$ and $\rho$, while passing to $\E^{\mathrm{co}}$ exchanges left and right adjoints and turns right mates into left mates of the transposed squares.

Let $\func(\Adj,\E)^{\mathrm{colax}}$ be the $(\infty,2)$-category of adjunctions in $\E$ and colax morphisms of adjunctions \cite[Def.~4.2]{haugseng2021laxmonads}, i.e.\ of colax transformations $\Adj \otimes \Delta^{1} \to \E$ for the Gray tensor product.
By \cite[Cor.~2.8.12]{abellangagnahaugseng2024straightening}, $\func(K,\E)$ is the wide locally full sub-$(\infty,2)$-category of $\func(K,\E)^{\mathrm{colax}}$ on the strong transformations, those whose naturality $2$-morphisms are invertible, naturally in $K$; so $\ell^{*}$ on $\func(\Adj,\E)$ is the restriction of $\ell^{*}\colon \func(\Adj,\E)^{\mathrm{colax}} \to \func(\Delta^{1},\E)^{\mathrm{colax}}$ to strong transformations.
By \cite[Cor.~4.8]{haugseng2021laxmonads}, the latter is an equivalence from $\func(\Adj,\E)^{\mathrm{colax}}$ onto the full sub-$(\infty,2)$-category of $\func(\Delta^{1},\E)$ on the left adjoints. Its proof uses \cite[Thm.~4.6]{haugseng2021laxmonads}, the lax analogue of \cref{thm:adjoints-in-functor-cats}, whose proof is completed in \cite[Cor.~5.2.10, Rem.~5.2.11]{abellangagnahaugseng2024straightening}.
So $\ell^{*}\colon \func(\Adj,\E) \to \func(\Delta^{1},\E)$ is the composite of a locally full inclusion, an equivalence and a full inclusion, hence a locally full inclusion, and its objects are the left adjoints.

A $1$-morphism of $\func(\Delta^{1},\E)$ between left adjoints $u$, $u'$ is a commutative square $\varphi$ as in (1). Under the equivalence of \cite[Cor.~4.8]{haugseng2021laxmonads} it corresponds to the colax morphism of adjunctions whose naturality square at $l$ is $\varphi$ and whose naturality square at $r$ is $\rmate{\varphi}$ \cite[Cor.~4.9]{haugseng2021laxmonads}. As $\Adj$ is generated by $l$ and $r$ under composition, and the naturality $2$-morphism of a colax transformation at a composite is the pasting of those at the factors, this colax morphism is strong if and only if $\rmate{\varphi}$ is invertible.
\end{proof}

\begin{rmk}
\label{rmk:mates-2cat}
For $2$-categories, \cref{lem:mates} is the classical fact that a morphism of adjunctions is determined by its left adjoint part, the right adjoint part being the mate \cite{kellystreet1974review}.
For $(\infty,2)$-categories, we know of no reference stating \cref{lem:mates} itself; the proof above deduces it from the lax version \cite[Cor.~4.8]{haugseng2021laxmonads}, in which $\ell^{*}$, resp.\ $\rho^{*}$, is an equivalence from adjunctions with colax, resp.\ lax, morphisms onto the full sub-$(\infty,2)$-category of $\func(\Delta^{1},\E)$ on the left, resp.\ right, adjoints. \cref{rmk:mates-alternative} gives a proof from \cref{thm:walking-adjunction,thm:adjoints-in-functor-cats} instead.
\end{rmk}

\subsubsection*{Adjunctions via lax pullbacks}

\begin{defi}
\label{def:lax-pullback}
A \emph{lax pullback} (comma object) of $f\colon X \to Y$ and $g\colon Z \to Y$ in $\E$ is an object $(f \downarrow g)$ with $1$-morphisms $p\colon (f \downarrow g) \to X$, $q\colon (f \downarrow g) \to Z$ and a $2$-morphism $\alpha\colon fp \Rightarrow gq$, such that for every $T \in \E$
\[
\Hom_{\E}\bigl(T,(f \downarrow g)\bigr) \longrightarrow (f_{\natural} \downarrow g_{\natural}) \ , \qquad h \mapsto (ph,qh,\alpha h) \ ,
\]
is an equivalence onto the comma $\infty$-category of triples $(u,v,fu \Rightarrow gv)$ with $u \in \Hom_{\E}(T,X)$ and $v \in \Hom_{\E}(T,Z)$.
\end{defi}

These equivalences are given by precomposition, hence natural in $T$. For $\E = \Cat$, $(f \downarrow g)$ is the comma $\infty$-category. We write $(f \downarrow Y) := (f \downarrow \mathrm{id}_{Y})$ and, for $g\colon Y \to X$, $(X \downarrow g) := (\mathrm{id}_{X} \downarrow g)$; both are spans from $X$ to $Y$.
For spans $S = (X \xleftarrow{p} E \xrightarrow{q} Y)$ and $S' = (X' \xleftarrow{p'} E' \xrightarrow{q'} Y')$ in $\E$ let
\[
\Hom(S,S') := \Hom(X,X') \times_{\Hom(E,X')} \Hom(E,E') \times_{\Hom(E,Y')} \Hom(Y,Y') \ ,
\]
the $\infty$-category of triples $(a,b,c)$ with $c\colon E \to E'$ and equivalences $p'c \simeq ap$, $q'c \simeq bq$. We call $c$ a \emph{map of spans over $(a,b)$}, and \emph{over $X \times Y$} if $X = X'$, $Y = Y'$ and $a$, $b$ are identities. The fibre of $\Hom(S,S')$ over $(a,b) \in \Hom(X,X') \times \Hom(Y,Y')$ is the $\infty$-category of maps of spans over $(a,b)$.

\begin{lem}
\label{lem:comma-fibrewise}
Let $F\colon A \to C$, $G\colon B \to C$, $F'\colon A \to C'$ and $G'\colon B \to C'$ be functors of $\infty$-categories. A functor $(F \downarrow G) \to (F' \downarrow G')$ over $A \times B$ is an equivalence if and only if it induces equivalences $\mapp(Fa,Gb) \to \mapp(F'a,G'b)$ on the fibres over all $(a,b)$.
\end{lem}

\begin{proof}
$(F \downarrow G) \to A \times B$ is the pullback of $\func(\Delta^{1},C) \to C \times C$ along $F \times G$, hence a bifibration \cite[Cor.~2.4.7.11]{htt}, with fibre $\mapp(Fa,Gb)$ over $(a,b)$. A functor over $A \times B$ between bifibrations that is an equivalence on fibres is an equivalence \cite[Prop.~2.4.7.6]{htt}.
\end{proof}

\begin{prop}[adjunctions via lax pullbacks]
\label{prop:comma-adjunctions}
Let $f\colon X \to Y$, $f'\colon X' \to Y'$, $g\colon Y \to X$ and $g'\colon Y' \to X'$ be $1$-morphisms of $\E$ whose lax pullbacks $(f \downarrow Y)$, $(f' \downarrow Y')$, $(X \downarrow g)$, $(X' \downarrow g')$ exist, and let
\[
i := (\mathrm{id}_{X},f,\mathrm{id}_{f})\colon X \to (f \downarrow Y) \ , \qquad j := (g,\mathrm{id}_{Y},\mathrm{id}_{g})\colon Y \to (X \downarrow g)
\]
be the generic sections.
\begin{enumerate}
\item Restriction along $i$, resp.\ $j$, gives equivalences over $\Hom(X,X') \times \Hom(Y,Y')$
\begin{gather*}
\Hom\bigl((f \downarrow Y),(f' \downarrow Y')\bigr) \simeq (f'_{\natural} \downarrow f^{\natural}) \ , \qquad
\Hom\bigl((f \downarrow Y),(X' \downarrow g')\bigr) \simeq (\mathrm{id} \downarrow g'_{\natural}f^{\natural}) \ , \\
\Hom\bigl((X \downarrow g),(X' \downarrow g')\bigr) \simeq (g^{\natural} \downarrow g'_{\natural}) \ ,
\end{gather*}
whose objects over $(a,b)$ are the lax squares $\sigma\colon f'a \Rightarrow bf$, resp.\ the $2$-morphisms $a \Rightarrow g'bf$, resp.\ the lax squares $\tau\colon ag \Rightarrow g'b$.
\item In particular, the $\infty$-category $\mapp_{X \times Y}((f \downarrow Y),(X \downarrow g))$ of maps of spans over $X \times Y$ is a space, and
\[
\mapp(\mathrm{id}_{X},gf) \xrightarrow{\ \simeq\ } \mapp_{X \times Y}\bigl((f \downarrow Y),(X \downarrow g)\bigr) \ , \qquad \eta \mapsto e_{\eta} := (p,q,g\alpha \circ \eta p) \ .
\]
The map $e_{\eta}$ is an equivalence in $\E$ if and only if $\eta$ is the unit of an adjunction $f \dashv g$. Hence $f$ is a left adjoint if and only if there are $g$ and an equivalence $(f \downarrow Y) \simeq (X \downarrow g)$ over $X \times Y$, and the space of such pairs $(g,e)$ is then contractible.
\item Let $f \dashv g$ and $f' \dashv g'$ with units $\eta$, $\eta'$, and let $e := e_{\eta}$, $e' := e_{\eta'}$. Conjugation $c \mapsto e'ce^{-1}$ and (1) give an equivalence over $\Hom(X,X') \times \Hom(Y,Y')$
\[
(f'_{\natural} \downarrow f^{\natural}) \xrightarrow{\ \simeq\ } (g^{\natural} \downarrow g'_{\natural})
\]
which sends a lax square $\sigma\colon f'a \Rightarrow bf$ to its right mate $\rmate{\sigma}\colon ag \Rightarrow g'b$.
\end{enumerate}
\end{prop}

\begin{proof}
(1) We have $pi = \mathrm{id}_{X}$, $qi = f$ and $\alpha i = \mathrm{id}_{f}$. The pair $(\mathrm{id}_{p},\alpha)$ is a $2$-morphism $ip \Rightarrow \mathrm{id}$ in $\Hom((f \downarrow Y),(f \downarrow Y)) \simeq (f_{\natural} \downarrow \mathrm{id})$, and it is the counit of an adjunction $i \dashv p$ with identity unit: the triangle identities read $p(\mathrm{id}_{p},\alpha) = \mathrm{id}_{p}$ and $(\mathrm{id}_{p},\alpha)i = (\mathrm{id},\alpha i) = \mathrm{id}_{i}$. Hence $p^{\natural} \dashv i^{\natural}$ (\cref{subsec:prelim-conventions}) with invertible unit, as $i^{\natural}p^{\natural} = (pi)^{\natural} = \mathrm{id}$, so that restriction along $i$ is an equivalence
\begin{equation}
\label{eq:comma-yoneda}
i^{\natural}\colon \mapp(hp,k) \xrightarrow{\ \simeq\ } \mapp(h,ki)
\end{equation}
for all $h\colon X \to T$ and $k\colon (f \downarrow Y) \to T$ \cite[Prop.~5.2.2.8]{htt}.
Write $E := (f \downarrow Y)$. By \cref{def:lax-pullback} for $T = E$, $\Hom(E,(f' \downarrow Y'))$ is the comma $\infty$-category $(f'_{\natural} \downarrow \mathrm{id})$ over $\Hom(E,X') \times \Hom(E,Y')$, so $\Hom((f \downarrow Y),(f' \downarrow Y'))$, its pullback along $p^{\natural} \times q^{\natural}$, is the comma $\infty$-category $(f'_{\natural}p^{\natural} \downarrow q^{\natural})$ of triples $(a,b,f'ap \Rightarrow bq)$. As $i^{\natural}p^{\natural} = \mathrm{id}$ and $i^{\natural}q^{\natural} = f^{\natural}$, restriction along $i$ is a functor $(f'_{\natural}p^{\natural} \downarrow q^{\natural}) \to (f'_{\natural} \downarrow f^{\natural})$ over $\Hom(X,X') \times \Hom(Y,Y')$, which on the fibres over $(a,b)$ is the equivalence \eqref{eq:comma-yoneda}. By \cref{lem:comma-fibrewise} it is an equivalence.
The second equivalence is the same with $(\mathrm{id} \downarrow g'_{\natural})$ in place of $(f'_{\natural} \downarrow \mathrm{id})$. The third is dual: $q \dashv j$ with identity counit and unit $(\beta,\mathrm{id}_{q})$, where $\beta\colon p' \Rightarrow gq'$ is the $2$-morphism of $(X \downarrow g)$, so $j^{\natural} \dashv q^{\natural}$ with invertible counit and restriction along $j$ is an equivalence $\mapp(k,hq) \simeq \mapp(kj,h)$.

(2) The first claim is the second equivalence of (1) on the fibres over $(\mathrm{id}_{X},\mathrm{id}_{Y})$, and the inverse sends $\eta$ to the map classified by $(p,q,g\alpha \circ \eta p)$, since $(g\alpha \circ \eta p)i = \eta$.
The map $e_{\eta}$ is an equivalence if and only if $\Hom(T,e_{\eta})$ is one for every $T$. By \cref{def:lax-pullback}, $\Hom(T,e_{\eta})$ is the functor $(f_{\natural} \downarrow \mathrm{id}) \to (\mathrm{id} \downarrow g_{\natural})$ over $\Hom(T,X) \times \Hom(T,Y)$ sending $(u,v,\varphi)$ to $(u,v,g\varphi \circ \eta u)$, so by \cref{lem:comma-fibrewise} $e_{\eta}$ is an equivalence if and only if
\begin{equation}
\label{eq:transposition}
\mapp(fu,v) \longrightarrow \mapp(u,gv) \ , \qquad \varphi \mapsto g\varphi \circ \eta u \ ,
\end{equation}
is an equivalence for all $T$, $u$, $v$. If $\eta$ is the unit of $f \dashv g$, then $f_{\natural} \dashv g_{\natural}$ with unit $\eta_{\natural}$ (\cref{subsec:prelim-conventions}), and \eqref{eq:transposition} is the transposition of this adjunction, an equivalence \cite[Prop.~5.2.2.8]{htt}. Conversely, suppose \eqref{eq:transposition} is an equivalence for all $T$, $u$, $v$. For $T = Y$, $u = g$, $v = \mathrm{id}_{Y}$, let $\varepsilon\colon fg \Rightarrow \mathrm{id}_{Y}$ be a preimage of $\mathrm{id}_{g}$, so that $g\varepsilon \circ \eta g = \mathrm{id}_{g}$ in $\htwo\E$. For $T = X$, $u = \mathrm{id}_{X}$, $v = f$, the map \eqref{eq:transposition} sends $\varepsilon f \circ f\eta$ to
\[
g\varepsilon f \circ gf\eta \circ \eta = g\varepsilon f \circ \eta gf \circ \eta = (g\varepsilon \circ \eta g)f \circ \eta = \eta \ ,
\]
the image of $\mathrm{id}_{f}$, so $\varepsilon f \circ f\eta = \mathrm{id}_{f}$ in $\htwo\E$: $\eta$ and $\varepsilon$ are the unit and counit of an adjunction $f \dashv g$ (\cref{def:adjunction}).
Finally, by the first two claims, the space of pairs $(g,e)$ is the space of pairs $(g,\eta)$ with $\eta$ the unit of an adjunction $f \dashv g$. Let $f \dashv f^{R}$ with unit $\eta_{0}$. The transposition $\mapp(f^{R},g) \simeq \mapp(\mathrm{id}_{X},gf)$, $\theta \mapsto \theta f \circ \eta_{0}$, of $(f^{R})^{\natural} \dashv f^{\natural}$ is natural in $g \in \Hom(Y,X)$ and identifies the invertible $\theta$ with the units, by uniqueness of adjoints in $\htwo\E$. So the space of pairs $(g,e)$ is the space of pairs $(g,\theta\colon f^{R} \simeq g)$, which is contractible.

(3) Let $\varepsilon$ be the counit of $f \dashv g$ and $p'$, $q'$, $\beta$ the structure of $(X \downarrow g)$. The map $d := (p',q',\varepsilon q' \circ f\beta)\colon (X \downarrow g) \to (f \downarrow Y)$ over $X \times Y$ is inverse to $e$: the composite $ed = (p',q',g\varepsilon q' \circ gf\beta \circ \eta p')$ restricts along $j$ to $(g,\mathrm{id}_{Y},g\varepsilon \circ \eta g) = (g,\mathrm{id}_{Y},\mathrm{id}_{g})$, as does $\mathrm{id}$, and restriction along $j$ is an equivalence on $\mapp_{X \times Y}((X \downarrow g),(X \downarrow g))$ by (1). By (1), the lax square $\sigma$ corresponds to $c_{\sigma} := (ap,bq,b\alpha \circ \sigma p)$, and $e'c_{\sigma}d$ restricts along $j$ to
\[
\bigl(ag,\,b,\,g'b\varepsilon \circ g'\sigma g \circ \eta' ag\bigr) \ ,
\]
whose $2$-morphism is $\rmate{\sigma}$ by \cref{constr:mates}.
\end{proof}

\begin{rmk}
\label{rmk:comma-literature}
When $\E$ is presented by an $\infty$-cosmos, \cref{prop:comma-adjunctions}(2) is \cite[Prop.~4.1.1]{riehlverity2022elements}, with the strict comma objects of the cosmos; for $\E = \Cat$ see \cite[Prop.~4.4.3]{riehlverity2015twocategory}. Since $\E$ is embedded in $\func(\E\opcat,\Cat)$ by (I3), the statement for $\E$ without lax pullbacks is (2) applied to $\Hom(-,f)$ and $\Hom(-,g)$.
\cref{prop:comma-adjunctions}(3) upgrades \cref{lem:mate-correspondence}(1) from spaces to $\infty$-categories. The equivalence $(f'_{\natural} \downarrow f^{\natural}) \simeq (g^{\natural} \downarrow g'_{\natural})$ itself only needs the adjunctions $f'_{\natural} \dashv g'_{\natural}$ and $g^{\natural} \dashv f^{\natural}$ of $\infty$-categories and \cref{lem:comma-fibrewise}; lax pullbacks realize it inside $\E$, as conjugation by $e$ and $e'$.
\end{rmk}

\begin{rmk}
\label{rmk:mates-alternative}
\cref{lem:mates} also follows from \cref{thm:walking-adjunction,thm:adjoints-in-functor-cats}. For $K \in \twoCat$, $\ell^{*}$ on $\mapp(K,-)$ is adjoint to $\ell^{*}\colon \mapp(\Adj,\func(K,\E)) \to \mapp(\Delta^{1},\func(K,\E))$, a monomorphism onto the left adjoints of $\func(K,\E)$ by \cref{thm:walking-adjunction}, and \cref{thm:adjoints-in-functor-cats} computes these for $K = \Delta^{0}$, $\Delta^{1}$, $\wtc$: for $K = \Delta^{1}$ they are the right Beck--Chevalley squares between left adjoints, and for $K = \wtc$ every $2$-morphism of $\func(\Delta^{1},\E)$ between two such squares is a left adjoint of $\func(\wtc,\E)$. So $\ell^{*}$ is a monomorphism and \eqref{eq:lff-square} is a pullback (\cref{subsec:prelim-conventions}), i.e.\ $\ell^{*}$ is a locally full inclusion.

When $\E$ has lax pullbacks, the two cases of \cref{thm:adjoints-in-functor-cats} used here follow from \cref{prop:comma-adjunctions}, granting that the hom-categories of $\func(K,\E)$ are the limits of the hom-categories of $\E$ indexed by $K$ \todo{Reference.}: then lax pullbacks in $\func(K,\E)$ exist and are computed pointwise, and a $1$-morphism of $\func(K,\E)$ is an equivalence if and only if its components are. For $K = \Delta^{1}$, let $L = (u,u',\varphi)\colon f_{0} \Rightarrow f_{1}$ be a $1$-morphism of $\func(\Delta^{1},\E)$ with $u$, $u'$ left adjoints. If $\rmate{\varphi}$ is invertible, then $R := (u^{R},u'^{R},\rmate{\varphi})\colon f_{1} \Rightarrow f_{0}$ is a $1$-morphism, and by \cref{prop:comma-adjunctions}(3) the pointwise equivalences $e_{u}$, $e_{u'}$ form an equivalence $(L \downarrow f_{1}) \simeq (f_{0} \downarrow R)$ over $f_{0} \times f_{1}$ in $\func(\Delta^{1},\E)$, so $L \dashv R$ by \cref{prop:comma-adjunctions}(2). Conversely, if $L \dashv R$, then $u \dashv R_{0}$ and $u' \dashv R_{1}$, and by \cref{prop:comma-adjunctions}(3) the naturality square of $R$ is $\rmate{\varphi}$, which is invertible because $R$ is a natural transformation. For $K = \wtc$ one argues in the same way, transporting the identification of the pastings of the two squares along the equivalence of $\infty$-categories in \cref{prop:comma-adjunctions}(3).
\end{rmk}

\subsubsection*{Reflections and coreflections}

\begin{defi}
\label{def:reflection}
An adjunction $l \dashv r$ in $\A$ is a \emph{reflection} if its counit $\varepsilon$ is invertible, and a \emph{coreflection} if its unit $\eta$ is invertible.
A $1$-morphism is \emph{reflective} if it is the right adjoint of a reflection, and \emph{coreflective} if it is the left adjoint of a coreflection.
\end{defi}

By \cref{thm:walking-adjunction} this is a property of $r$, resp.\ $l$. As $r_{\natural}$ and $l_{\natural}$ on $\Hom_{\A}(t,-)$ form an adjunction with counit $\varepsilon_{\natural}$ and unit $\eta_{\natural}$, a right adjoint $r$ is reflective if and only if every $r_{\natural}$ is fully faithful, and a left adjoint $l$ is coreflective if and only if every $l_{\natural}$ is fully faithful.

Let $c\colon \wtc \to \Delta^{1}$ collapse the $2$-morphism, and let $\varepsilon,\eta\colon \wtc \to \Adj$ classify the counit $lr \Rightarrow \mathrm{id}_{+}$ and the unit $\mathrm{id}_{-} \Rightarrow rl$. The \emph{walking reflection} and \emph{walking coreflection} are the pushouts in $\twoCat$
\[
\Refl := \Adj \sqcup_{\wtc,\varepsilon} \Delta^{1} \ , \qquad \Corefl := \Adj \sqcup_{\wtc,\eta} \Delta^{1}
\]
along $c$. We write $\ell$, $\rho$ also for the composites $\Delta^{1} \to \Adj \to \Refl$, $\Corefl$.

\begin{lem}[mates for reflections]
\label{lem:mates-refl}
Let $\E$ be an $(\infty,2)$-category.
\begin{enumerate}
\item $\func(\Refl,\E) \to \func(\Adj,\E)$ is fully faithful and exhibits $\func(\Refl,\E)$ as the full sub-$(\infty,2)$-category on the reflections; likewise $\func(\Corefl,\E)$ is the full sub-$(\infty,2)$-category on the coreflections.
\item Hence for $X \in \{\Refl,\Corefl\}$, $\ell^{*}$ and $\rho^{*}\colon \func(X,\E) \to \func(\Delta^{1},\E)$ are locally full inclusions with the $1$-morphisms described in \cref{lem:mates}. The objects in the image of $\ell^{*}$, resp.\ $\rho^{*}$, are the left, resp.\ right, adjoints of reflections for $X = \Refl$, and of coreflections for $X = \Corefl$.
\end{enumerate}
\end{lem}

\begin{proof}
(1) $\mapp(\wtc,\B)$ is the space of objects $x,y \in \B$ with a morphism of $\Hom_{\B}(x,y)$, and $c^{*}$ sends $f$ to $\mathrm{id}_{f}$. As $C^{\simeq} \to \func(\Delta^{1},C)^{\simeq}$ is a monomorphism onto the equivalences for every $\infty$-category $C$, $c^{*}$ is a monomorphism onto the invertible $2$-morphisms.
Applying this to $\B = \func(K,\E)$, the map
\[
\mapp\bigl(K,\func(\Refl,\E)\bigr) \simeq \mapp\bigl(\Refl,\func(K,\E)\bigr) \longrightarrow \mapp\bigl(\Adj,\func(K,\E)\bigr) \simeq \mapp\bigl(K,\func(\Adj,\E)\bigr)
\]
is a monomorphism onto the adjunctions in $\func(K,\E)$ with invertible counit. A $2$-morphism of $\func(K,\E)$ is invertible if and only if its components at the objects of $K$ are.\todo{Reference.}
So for $K = \Delta^{0}$, $\Delta^{1}$, $\wtc$ the image consists of all $K \to \func(\Adj,\E)$ whose objects are reflections. The case of $\Corefl$ is the same.
(2) follows from (1) and \cref{lem:mates}.
\end{proof}
